\section{化学平衡・反応}

\subsection{平成25年度 問13}
\label{q:h25-reg-applied-13}

\noindent\textbf{原典：} sources/H25/17\_applied\_science.pdf、PDF物理6頁。

\textbf{公式正答：} 4（single）。

\par\textit{原典本文の校正と独自解説は未作成。分類のみの登録項目。}

\section{大気・海洋・物質循環}

\subsection{平成26年度 問24}
\label{q:h26-reg-applied-24}

\noindent\textbf{原典：} sources/H26/17\_applied\_science.pdf、PDF物理10頁。

\textbf{公式正答：} 2（single）。

\par\textit{原典本文の校正と独自解説は未作成。分類のみの登録項目。}

\subsection{平成27年度 問23}
\label{q:h27-reg-applied-23}

\noindent\textbf{原典：} sources/H27/17\_applied\_science.pdf、PDF物理9頁。

\textbf{公式正答：} 3（single）。

\par\textit{原典本文の校正と独自解説は未作成。分類のみの登録項目。}

\section{フッ化水素の物性と酸解離}

\subsection{令和元年度 問9}
\label{q:r01-reg-applied-09}

\noindent\textbf{原典：} sources/R01/17\_applied\_science.pdf、PDF物理3頁。

\textbf{公式正答：} 5（single）。

\par\textbf{検討未確定：} 純液体HFの誘電率と単独HFの金・白金非反応性について条件付き一次資料の全文確認が残る。発煙も濃度条件を要する。初回5と公式5の一致だけで全肢校正完了としない。

\paragraph{問題の要旨} フッ化水素の物性・ガラスへの反応・危険性と、酸解離の強さを区別する。

\Needspace{4\baselineskip}\paragraph{確認した条件・記述（独自要約）}\mbox{}\par

\noindent\textbf{確認範囲：} 問題幹、否定、全5肢、数式、数値、単位を原典画像で確認した。\par

\noindent\textbf{成立条件：} 他のハロゲン化水素との比較。純液体の誘電性と水溶液の酸解離は別。単独のHFと他の酸化剤を併用する系を区別する。\par

\paragraph{問われる核心} フッ化水素の物性・ガラスへの反応・危険性と、酸解離の強さを区別する。

\paragraph{必要知識} 他のハロゲン化水素との比較。純液体の誘電性と水溶液の酸解離は別。単独のHFと他の酸化剤を併用する系を区別する。

\paragraph{考え方・正答を選ぶ理由（査読待ち初稿）} HFは強い水素結合を形成し、他のハロゲン化水素より高い沸点を持つ。ガラスのケイ酸成分と反応する。高濃度水溶液は発煙し、弱酸であっても深刻な有害性を持つ。5は有害性を強酸であることに結び付けるため誤りである。

\Needspace{4\baselineskip}\paragraph{選択肢の検討（査読待ち初稿）}\mbox{}\par

\noindent\textbf{1：適切} 他のハロゲン化水素より沸点が高い。分子間水素結合による会合が寄与する。\par

\noindent\textbf{2：解釈未解決} 液体HFは大きな誘電率を持つ。定性的な候補として初回は適切としたが、温度条件付きの定量一次資料の照合が残る。\par

\noindent\textbf{3：条件付きで適切} ガラスと反応するが金・白金とは反応しない。ガラス腐食はNIOSHで確認できる。金・白金の非反応性を他の酸化剤との混合条件まで一般化しない。\par

\noindent\textbf{4：条件付きで適切} 水溶液は空気中で発煙する。NIOSHは濃度40パーセント超の発煙を説明する。濃度を指定しない全水溶液への全称とは読まない。\par

\noindent\textbf{5：不適切} 強い毒性を持つ強酸である。毒性・腐食性は強いが、通常の希薄水溶液における解離は弱酸の分類である。\par

\paragraph{再利用する考え方} 反応速度と平衡位置、全圧と分圧、濃度と活量を分ける。式の適用条件を確認し、原典の不正確な記述を科学的に正しい文へ置き換えない。

\paragraph{共通理論} \hyperref[theory:chem_eq:acid_strength]{第\ref*{theory:chem_eq:acid_strength}節}

\par 参考：\cite{purdue_acid_strength,niosh_hf}。

\par 難易度 2（校準初稿）、優先度未評価。

\section{未確定}

\subsection{平成23年度 問11}
\label{q:h23-reg-applied-11}

\noindent\textbf{原典：} sources/H23/17\_applied\_science.pdf、PDF物理4頁。

\textbf{公式正答：} 4（single）。

\par\textit{原典本文の校正と独自解説は未作成。分類のみの登録項目。}

\subsection{平成23年度 問13}
\label{q:h23-reg-applied-13}

\noindent\textbf{原典：} sources/H23/17\_applied\_science.pdf、PDF物理5頁。

\textbf{公式正答：} 5（single）。

\par\textit{原典本文の校正と独自解説は未作成。分類のみの登録項目。}

\subsection{平成23年度 問17}
\label{q:h23-reg-applied-17}

\noindent\textbf{原典：} sources/H23/17\_applied\_science.pdf、PDF物理6頁。

\textbf{公式正答：} 1（single）。

\par\textit{原典本文の校正と独自解説は未作成。分類のみの登録項目。}

\subsection{平成24年度 問12}
\label{q:h24-reg-applied-12}

\noindent\textbf{原典：} sources/H24/17\_applied\_science.pdf、PDF物理4頁。

\textbf{公式正答：} 4（single）。

\par\textit{原典本文の校正と独自解説は未作成。分類のみの登録項目。}

\subsection{平成26年度 問12}
\label{q:h26-reg-applied-12}

\noindent\textbf{原典：} sources/H26/17\_applied\_science.pdf、PDF物理5頁。

\textbf{公式正答：} 1（single）。

\par\textit{原典本文の校正と独自解説は未作成。分類のみの登録項目。}

\subsection{平成27年度 問10}
\label{q:h27-reg-applied-10}

\noindent\textbf{原典：} sources/H27/17\_applied\_science.pdf、PDF物理4頁。

\textbf{公式正答：} 4（single）。

\par\textit{原典本文の校正と独自解説は未作成。分類のみの登録項目。}

\subsection{平成28年度 問11}
\label{q:h28-reg-applied-11}

\noindent\textbf{原典：} sources/H28/17\_applied\_science.pdf、PDF物理4頁。

\textbf{公式正答：} 5（single）。

\par\textit{原典本文の校正と独自解説は未作成。分類のみの登録項目。}

\subsection{平成29年度 問10}
\label{q:h29-reg-applied-10}

\noindent\textbf{原典：} sources/H29/17\_applied\_science.pdf、PDF物理3頁。

\textbf{公式正答：} 4（single）。

\par\textit{原典本文の校正と独自解説は未作成。分類のみの登録項目。}

\subsection{平成29年度 問8}
\label{q:h29-reg-applied-08}

\noindent\textbf{原典：} sources/H29/17\_applied\_science.pdf、PDF物理3頁。

\textbf{公式正答：} 5（single）。

\par\textit{原典本文の校正と独自解説は未作成。分類のみの登録項目。}

\subsection{平成30年度 問10}
\label{q:h30-reg-applied-10}

\noindent\textbf{原典：} sources/H30/17\_applied\_science.pdf、PDF物理3頁。

\textbf{公式正答：} 3（single）。

\par\textit{原典本文の校正と独自解説は未作成。分類のみの登録項目。}

\subsection{令和元年度再試験 問16}
\label{q:r01-ret-applied-16}

\noindent\textbf{原典：} sources/R01\_retest/17\_applied\_science.pdf、PDF物理6頁。

\textbf{公式正答：} 2（single）。

\par\textit{原典本文の校正と独自解説は未作成。分類のみの登録項目。}

\subsection{令和元年度再試験 問9}
\label{q:r01-ret-applied-09}

\noindent\textbf{原典：} sources/R01\_retest/17\_applied\_science.pdf、PDF物理4頁。

\textbf{公式正答：} 3（single）。

\par\textit{原典本文の校正と独自解説は未作成。分類のみの登録項目。}

\subsection{令和2年度 問14}
\label{q:r02-reg-applied-14}

\noindent\textbf{原典：} sources/R02/17\_applied\_science.pdf、PDF物理5頁。

\textbf{公式正答：} 4（single）。

\par\textit{原典本文の校正と独自解説は未作成。分類のみの登録項目。}

\subsection{令和2年度 問16}
\label{q:r02-reg-applied-16}

\noindent\textbf{原典：} sources/R02/17\_applied\_science.pdf、PDF物理6頁。

\textbf{公式正答：} 5（single）。

\par\textit{原典本文の校正と独自解説は未作成。分類のみの登録項目。}

\subsection{令和3年度 問13}
\label{q:r03-reg-applied-13}

\noindent\textbf{原典：} sources/R03/17\_applied\_science.pdf、PDF物理5頁。

\textbf{公式正答：} 1（single）。

\par\textit{原典本文の校正と独自解説は未作成。分類のみの登録項目。}

\subsection{令和4年度 問13}
\label{q:r04-reg-applied-13}

\noindent\textbf{原典：} sources/R04/17\_applied\_science.pdf、PDF物理4頁。

\textbf{公式正答：} 1（single）。

\par\textit{原典本文の校正と独自解説は未作成。分類のみの登録項目。}

\subsection{令和5年度 問12}
\label{q:r05-reg-applied-12}

\noindent\textbf{原典：} sources/R05/17\_applied\_science.pdf、PDF物理5頁。

\textbf{公式正答：} 5（single）。

\par\textit{原典本文の校正と独自解説は未作成。分類のみの登録項目。}

\subsection{令和6年度 問10}
\label{q:r06-reg-applied-10}

\noindent\textbf{原典：} sources/R06/17\_applied\_science.pdf、PDF物理3頁。

\textbf{公式正答：} 1（single）。

\par\textit{原典本文の校正と独自解説は未作成。分類のみの登録項目。}

\subsection{令和6年度 問13}
\label{q:r06-reg-applied-13}

\noindent\textbf{原典：} sources/R06/17\_applied\_science.pdf、PDF物理4頁。

\textbf{公式正答：} 5（single）。

\par\textit{原典本文の校正と独自解説は未作成。分類のみの登録項目。}

\subsection{令和6年度 問9}
\label{q:r06-reg-applied-09}

\noindent\textbf{原典：} sources/R06/17\_applied\_science.pdf、PDF物理3頁。

\textbf{公式正答：} 5（single）。

\par\textit{原典本文の校正と独自解説は未作成。分類のみの登録項目。}

\subsection{令和7年度 問15}
\label{q:r07-reg-applied-15}

\noindent\textbf{原典：} sources/R07/17\_applied\_science.pdf、PDF物理5頁。

\textbf{公式正答：} 2（single）。

\par\textit{原典本文の校正と独自解説は未作成。分類のみの登録項目。}

\subsection{令和7年度 問16}
\label{q:r07-reg-applied-16}

\noindent\textbf{原典：} sources/R07/17\_applied\_science.pdf、PDF物理6頁。

\textbf{公式正答：} 5（single）。

\par\textit{原典本文の校正と独自解説は未作成。分類のみの登録項目。}

\subsection{令和7年度 問17}
\label{q:r07-reg-applied-17}

\noindent\textbf{原典：} sources/R07/17\_applied\_science.pdf、PDF物理6頁。

\textbf{公式正答：} 3（single）。

\par\textit{原典本文の校正と独自解説は未作成。分類のみの登録項目。}
